% Part: turing-machines
% Chapter: undecidability
% Section: introduction

\documentclass[../../../include/open-logic-section]{subfiles}

\begin{document}

\olfileid{tur}{und}{int}
\olsection{પરિચય}

પ્રથમ નજરે એવું સ્પષ્ટ લાગશે કે કેટલાક વિધેયો, અંકગણિતીય
વિધેયો પણ, અસંગણનીય છે. આવાં વિધેયો એટલાં બધાં છે
અને તેમનું વર્તન એટલું
જટિલ છે. દાખલા તરીકે, કિરણોત્સર્ગી કણોના ક્ષય પરથી, અથવા
બીજા અવ્યવસ્થિત કે યાદૃચ્છિક વર્તન પરથી વ્યાખ્યાયિત વિધેયો લો.
ધારો કે આપણે કોઈ આપેલી ક્ષણથી એક-સેકંડના સમયગાળાની ગણતરી
શરૂ કરીએ અને તે પ્રારંભિક ક્ષણ પછીના $n$મા એક-સેકંડના
સમયગાળામાં બ્રહ્માંડના ક્ષય પામતા કણોની સંખ્યા $f(n)$ તરીકે
વ્યાખ્યાયિત કરીએ. આ વિધેય કદાચ એવું લાગે કે તેને કદી
ગણવાની આશા રાખી શકાય નહીં.

પરંતુ આવા વિધેયો કેવી રીતે ગણવા તેની કલ્પના ન થઈ શકે એ એક
વાત છે; તેઓ અસંગણનીય છે તે ખરેખર સાબિત કરવું બીજી વાત છે.
હકીકતમાં, અત્યંત જટિલ લાગતાં વિધેયો પણ અમૂર્ત અર્થમાં
સંગણનીય હોઈ શકે. ઉદાહરણ તરીકે, ધારો કે બ્રહ્માંડ સમયની
દૃષ્ટિએ સાન્ત છે---કેટલાક બ્રહ્માંડવિજ્ઞાનના સિદ્ધાંતોની
આગાહી મુજબ, દૂરના ભવિષ્યમાં કોઈ દિવસે તે સંકોચાઈને એક જ
બિંદુ બની જશે. તો એ પ્રારંભિક ક્ષણ પછી $f(n)$ વ્યાખ્યાયિત
હોય એવા સેકંડોની સંખ્યા સાન્ત (જોકે અતિ વિશાળ) હશે.
માત્ર સાન્ત સંખ્યામાં ઇનપુટ માટે વ્યાખ્યાયિત કોઈ પણ વિધેય
સંગણનીય છે: આપણે તેના આઉટપુટ એક મોટા કોષ્ટકમાં લખી શકીએ,
અથવા તેમને અત્યંત મોટા ટ્યુરિંગ મશીનના અવસ્થા-આલેખમાં
સંકેતિત કરી શકીએ.

ઘણી વાર આપણને એવા વિધેયોના ખાસ કિસ્સાઓમાં રસ હોય છે
જેનાં મૂલ્યો હા/ના પ્રશ્નોના જવાબ આપે છે. ઉદાહરણ તરીકે,
``$n$ અવિભાજ્ય સંખ્યા છે?'' એ પ્રશ્નની સાથે નીચેનું વિધેય જોડાય છે:
\[
\fn{isprime}(n) = \begin{cases}
  1 & \text{જો $n$ અવિભાજ્ય હોય}\\
  0 & \text{અન્યથા.}
  \end{cases}
\]
સંબંધિત $1/0$-મૂલ્યવાળું વિધેય અસરકારક રીતે સંગણનીય હોય,
તો આપણે કહીએ કે હા/ના પ્રશ્નનો \emph{અસરકારક રીતે નિર્ણય}
કરી શકાય છે.

કેટલાંક વિધેયો અસરકારક રીતે ગણી શકાતાં નથી, અથવા કેટલીક
સમસ્યાઓનો અસરકારક રીતે નિર્ણય કરી શકાતો નથી, તે ગણિતીય
રીતે સાબિત કરવા સંગણનનું કોઈ એક ચોક્કસ નિદર્શ નક્કી કરવું
અને તે કેટલાંક વિધેયો ગણી કે કેટલીક સમસ્યાઓનો નિર્ણય
કરી શકતું નથી તે બતાવવું આવશ્યક છે. દાખલા તરીકે, દરેક
વિધેય ટ્યુરિંગ મશીન વડે ગણી શકાતું નથી અને દરેક સમસ્યાનો
ટ્યુરિંગ મશીન વડે નિર્ણય કરી શકાતો નથી તે આપણે બતાવી
શકીએ. પછી ચર્ચ--ટ્યુરિંગ માન્યતાનો આધાર લઈને તારવી
શકીએ કે માત્ર ટ્યુરિંગ મશીનો જ દરેક વિધેય ગણવા માટે
અસમર્થ નથી; ટ્યુરિંગ મશીન વડે ન ગણી શકાય તેવા વિધેયો
કોઈ અસરકારક પ્રક્રિયા વડે પણ ગણી શકાતા નથી.

આવા નકારાત્મક પરિણામોની સાબિતીમાં મુખ્ય બાબત એ છે કે
ટ્યુરિંગ મશીનોને જ સંખ્યાઓ આપી શકાય છે. એ માટેનો સૌથી
સરળ રસ્તો તેમની પરિગણના કરવાનો છે: મશીનો અને તેમના
પ્રોગ્રામ લખવાની કોઈ ચોક્કસ રીત નક્કી કરી, પછી તેમને
ક્રમબદ્ધ રીતે યાદીમાં મૂકી શકાય. આમ કરી શકાય છે એ
જોયા પછી, ટ્યુરિંગ મશીન વડે અસંગણનીય વિધેયોનું અસ્તિત્વ
ગણાંકના સરળ વિચારોથી ફલિત થાય છે: $\Nat$થી~$\Nat$માં
જતાં વિધેયોનો ગણ (હકીકતમાં, $\Nat$થી~$\{0, 1\}$માં
જતાં વિધેયોનો ગણ પણ) !!{nonenumerable} છે; પરંતુ બધાં
ટ્યુરિંગ મશીનોની પરિગણના કરી શકાય છે, તેથી ટ્યુરિંગ-સંગણનીય
વિધેયોનો ગણ માત્ર !!{denumerable} છે.

આપણે \emph{ચોક્કસ} વિધેયો અને સમસ્યાઓ પણ વ્યાખ્યાયિત
કરી શકીએ, અને અનુક્રમે તેમનું અસંગણનીય તથા અનિર્ણેય
હોવું સાબિત કરી શકીએ. આવી એક સમસ્યા કહેવાતી
\emph{અટકવાની સમસ્યા} છે. ટ્યુરિંગ મશીનની સૂચનાઓની
યાદી આપીને તેનું સાન્ત વર્ણન કરી શકાય છે. ટ્યુરિંગ
મશીનનું એવું વર્ણન, એટલે ટ્યુરિંગ મશીનનો પ્રોગ્રામ,
અલબત્ત બીજા ટ્યુરિંગ મશીનને ઇનપુટ તરીકે આપી શકાય.
તેથી આપણે એવા ટ્યુરિંગ મશીનો વિચારીએ જે બીજાં
ટ્યુરિંગ મશીનો વિશેના પ્રશ્નોનો નિર્ણય કરે. ખાસ
રસપ્રદ પ્રશ્ન આ છે: ``આપેલું ટ્યુરિંગ મશીન ઇનપુટ~$n$
પર શરૂ કરતાં છેવટે અટકશે?'' આ પ્રશ્નનો નિર્ણય
કરી શકતું ટ્યુરિંગ મશીન હોય તો સારું: તેને ગુણવત્તા
તપાસતું ટ્યુરિંગ મશીન માનો, જે બીજા મશીનો અનંત
ચક્રમાં ફસાઈ ન જાય તેની ખાતરી કરે. ટ્યુરિંગે સાબિત
કરેલી રસપ્રદ હકીકત એ છે કે આવું મશીન હોઈ શકે નહીં.
ટ્યુરિંગ મશીન~$M$ના વર્ણન અને કોઈ સંખ્યા~$n$નો
ઇનપુટ આપતાં હંમેશાં અટકે અને, $M$ ઇનપુટ~$n$ પર
અટકે કે ન અટકે તે મુજબ અનુક્રમે $1$ કે $0$ આપે,
એવું કોઈ એક ટ્યુરિંગ મશીન હોઈ શકે નહીં.

ચોક્કસ અનિર્ણેય સમસ્યાઓનાં ઉદાહરણો મળ્યા પછી બીજી
સમસ્યાઓ પણ અનિર્ણેય છે તે બતાવવા તેમનો ઉપયોગ
કરી શકીએ. દાખલા તરીકે, પ્રસિદ્ધ અનિર્ણેય પ્રશ્ન છે:
``પ્રથમ-ક્રમનું !!{formula}~$!A$ પ્રામાણ્ય ધરાવે છે?''
ઇનપુટ તરીકે પ્રથમ-ક્રમનું !!{formula}~$!A$ આપવામાં આવે
ત્યારે, $!A$ પ્રામાણ્ય ધરાવે છે કે નહીં તે મુજબ $1$ કે
$0$ આઉટપુટ સાથે અટકશે તેની ખાતરી હોય એવું કોઈ ટ્યુરિંગ
મશીન નથી. ઐતિહાસિક રીતે, આ સમસ્યાને અસરકારક રીતે
ઉકેલવાની પ્રક્રિયા શોધવાના પ્રશ્નને સાદી રીતે ``નિર્ણય
સમસ્યા'' કહેવાતો; તેથી આપણે કહીએ છીએ કે નિર્ણય સમસ્યા
ઉકેલી શકાતી નથી. ટ્યુરિંગ અને ચર્ચે આશરે એક જ સમયે સ્વતંત્ર
રીતે આ પરિણામ સાબિત કર્યું હતું, તેથી તેને ચર્ચ--ટ્યુરિંગ
પ્રમેય પણ કહે છે.

\end{document}
